\documentclass[11pt]{article}
\usepackage[a4paper,margin=26mm]{geometry}
\usepackage[T1]{fontenc}
\usepackage{lmodern}
\usepackage{amsmath,amssymb,amsthm}
\usepackage{microtype}
\usepackage{booktabs}
\usepackage{xurl}
\usepackage{hyperref}
\hypersetup{colorlinks=true,linkcolor=blue,urlcolor=blue,citecolor=blue,
 pdftitle={Critical graphs with the fewest cycles of length 0 modulo k},
 pdfauthor={Horace Dong}}
\newtheorem{theorem}{Theorem}
\newtheorem{lemma}{Lemma}
\newtheorem{proposition}{Proposition}
\newtheorem{corollary}{Corollary}
\theoremstyle{remark}
\newtheorem{remark}{Remark}
\newcommand{\repoversion}{1.22.0}
\setlength{\emergencystretch}{1.5em}

\title{Critical graphs with the fewest cycles of length $0$ modulo $k$}
\author{Horace Dong\\ \small Independent researcher\\
\small\texttt{maxwelldhx@gmail.com}\\
\small ORCID: \href{https://orcid.org/0009-0002-7554-7548}{0009-0002-7554-7548}}
\date{18 November 2026}
\begin{document}
\maketitle

\begin{abstract}
Choi and Sieniawski proved that for every $k\ge3$, every graph with no proper
$k$-colouring has at least $(k+1)!/(2k)$ cycles of length $0$ modulo~$k$, a
bound attained by $K_{k+1}$, and asked whether, for $k\ge4$, $K_{k+1}$ is the
only $(k+1)$-critical graph attaining it. We show that the answer is yes for
every $k\ge4$. The proof uses the counting method of Choi and Sieniawski. For
$k=4$ and $k=5$ it uses a construction of Kim and Picollelli and a case
analysis of degrees. For $k\ge6$ we show that a critical graph attaining the
bound, other than $K_{k+1}$, has minimum degree at least $2k-1$, and then use
a theorem of Gao, Huo, Liu and Ma on paths whose lengths form an arithmetic
progression to find a cycle whose length is a multiple of $k$ larger than $k$,
which the equality case forbids. The result is formally verified in Lean~4,
including the theorem of Gao, Huo, Liu and Ma.
The proofs were found and formalized by AI agents working under the author's
direction; their role is described at the end of the paper.
\end{abstract}

\section{Introduction}

All graphs are finite and simple. A cycle is identified with its edge set, so
neither a starting vertex nor a direction is chosen. For $k\ge3$ and a graph $G$
let $N_k(G)$ be the number of cycles of $G$ whose length is divisible by~$k$.
Brewster, McGuinness, Moore and Noel~\cite{bmmn} conjectured that $K_{k+1}$
minimises $N_k$ among graphs with no proper $k$-colouring. Kim and
Picollelli~\cite{kp} proved this for $k=3$, and Choi and
Sieniawski~\cite{cs} proved it for every $k\ge3$:

\begin{theorem}[{Choi and Sieniawski~\cite[Theorem~3]{cs}}]\label{thm:cs}
For every $k\ge3$, every graph with no proper $k$-colouring has at least
$\frac{(k+1)!}{2k}$ cycles of length $0$ modulo~$k$.
\end{theorem}

The bound is attained by $K_{k+1}$: its only cycles of length divisible by $k$
are its $k$-cycles, and there are $(k+1)\cdot\frac{(k-1)!}{2}=\frac{(k+1)!}{2k}$
of them. A graph is \emph{$(k+1)$-critical} if its chromatic number is $k+1$ and
every proper subgraph (including those obtained by deleting a single edge) has a
proper $k$-colouring. Kim and Picollelli showed that $K_4$ is the only
$4$-critical graph with exactly four cycles of length $0$ modulo~$3$. Choi and
Sieniawski ask~\cite[Question~8]{cs}:
\begin{quote}
Let $k\ge4$. Is $K_{k+1}$ the unique $(k+1)$-critical graph with exactly
$\frac{(k+1)!}{2k}$ cycles of length $0$ modulo $k$?
\end{quote}

\begin{theorem}\label{thm:main}
Let $k\ge4$. If $G$ is a $(k+1)$-critical graph with exactly
$\frac{(k+1)!}{2k}$ cycles of length $0$ modulo $k$, then $G\cong K_{k+1}$.
\end{theorem}

So the answer to the question is yes. A first version of this paper
(version~1.21.0 of the repository~\cite{repo}) proved Theorem~\ref{thm:main}
for $4\le k\le7$, with a separate case analysis for each~$k$. The present
version keeps the arguments for $k=4$ and $k=5$ and replaces those for $k=6$
and $k=7$ by one argument for all $k\ge6$ (Sections~\ref{sec:six}
and~\ref{sec:long}).

For graphs of minimum degree $k$ we use a construction of Kim and
Picollelli~\cite[Theorem~2.4 and Corollary~2.5]{kp}: for a $(k+1)$-critical
graph $G$ with $\delta(G)=k$ they construct an injection $f$ from the cycles of
$K_{k+1}$ of length at most $k$ into the cycles of $G$ such that the length of
$f(C)$ is a multiple of the length of $C$, and deduce the lower bound
$(k+1)!/(2k)$ for such graphs. They do not consider the equality case of this
bound for $k\ge4$. For $k\ge6$ we use a theorem of Gao, Huo, Liu and
Ma~\cite[Theorem~3.1]{ghlm} on paths with prescribed ends whose lengths form an
arithmetic progression with difference $1$ or $2$, which in turn uses a theorem of
Fan~\cite{fan}.

Searches of the literature on 13, 16, 17 and 18 November 2026 found no answer to the
question for any $k\ge4$ other than version~1 of this paper; arXiv:2611.13580
had only its first version and no citing works. A public repository of
AI-written notes treats the lower bound for special classes of graphs, such as
line graphs, but not this question for general critical graphs. The coverage
of these searches is described in the repository~\cite{repo}.

\section{Good cycles and the saturation identity}\label{sec:good}

Fix $k\ge4$ and write $b=\frac{(k-1)!}{2}$ and $B=(k+1)b=\frac{(k+1)!}{2k}$.
We use the method of~\cite{cs}. Let a total order of $V(G)$ be given. An edge of
a directed walk is a \emph{descent} if it goes from a later to an earlier
vertex. A cycle of length $ak$ with a chosen cyclic direction is \emph{good} if
it has at most $a$ descents. The two directions of a cycle of length $ak$ have
$ak$ descents together, and $ak>2a$, so at most one direction is good; a cycle
is good if one of its directions is.

Let $v$ be a vertex such that $G-v$ has a proper $k$-colouring $\varphi$ with
colours $1,\dots,k$. For a permutation $\pi$ of the colours, order $V(G)-v$ by
the colours $\pi(\varphi(\cdot))$ (ties broken by a fixed order) and put $v$
last; call this the order induced by $(\varphi,\pi)$.

\begin{lemma}[{\cite[Lemmas~5 and~6]{cs}}]\label{lem:om}
Suppose $G$ has no proper $k$-colouring.
\begin{enumerate}
\item[(O)] For every $\pi$, the order induced by $(\varphi,\pi)$ has at least
$k+1$ good cycles.
\item[(M)] Every cycle of length divisible by $k$ is good for at most $2k$ of the
$k!$ permutations~$\pi$.
\end{enumerate}
\end{lemma}

Counting pairs $(\pi,C)$ with $C$ good for $\pi$ gives
$k!(k+1)\le 2k\,N_k(G)$ whenever some $G-v$ is $k$-colourable; applied to a
minimal subgraph that is not $k$-colourable, this gives Theorem~\ref{thm:cs}. If $N_k(G)=B$, both
inequalities are equalities.

\begin{corollary}[saturation]\label{cor:sat}
If $G$ has no proper $k$-colouring and $N_k(G)=B$, then for every vertex $v$,
every proper $k$-colouring $\varphi$ of $G-v$ and every cycle $C$ of length
divisible by $k$, the cycle $C$ is good for exactly $2k$ permutations~$\pi$.
\end{corollary}

We recall the part of the proof of (M) that we use. Let $C$ have length $ak$ and
pass through $v$, and write a direction of $C$ as
$u_1u_2\cdots u_{ak-1}vu_1$. The edge $vu_1$ is a descent, so the path
$u_1\cdots u_{ak-1}$ has at most $a-1$ descents. Its last $(a-1)k$ edges split
into $a-1$ blocks of $k$ consecutive edges, and each block contains a descent
(otherwise $k+1$ vertices would have strictly increasing colours). Hence the
first $k-2$ edges are ascents: the colours
$\alpha_1=\varphi(u_1),\dots,\alpha_{k-1}=\varphi(u_{k-1})$ are distinct and
$\pi$ orders them increasingly. The permutations for which this direction is
good are therefore among the $k$ permutations obtained by inserting the
remaining colour $\beta$ into the increasing sequence
$\alpha_1<\dots<\alpha_{k-1}$. If $C$ avoids $v$, a good direction has exactly
one descent in every block of $k$ consecutive edges, so the colours around $C$
repeat a fixed cyclic order, and again at most $k$ permutations make it good.

\section{Consequences of saturation}\label{sec:general}

Throughout this section $k\ge4$ and $G$ is $(k+1)$-critical with $N_k(G)=B$. Since $G$ is
critical, $G-v$ is $k$-colourable for every $v$, and every vertex has degree at
least~$k$.

\begin{lemma}\label{lem:short}
Every cycle of $G$ whose length is divisible by $k$ has length $k$.
\end{lemma}

\begin{proof}
Suppose $C$ has length $ak$ with $a\ge2$. Pick $v\in V(C)$ and a proper
$k$-colouring $\varphi$ of $G-v$, and fix a direction
$u_1\cdots u_{ak-1}vu_1$ of $C$; keep the notation of Section~\ref{sec:good}.
We show that this direction is good for at most $k-1$ permutations. Then $C$ is
good for at most $2k-2$ permutations, contradicting
Corollary~\ref{cor:sat}.

Suppose the direction is good for all $k$ candidate permutations. Take first the
permutation with $\beta<\alpha_1<\dots<\alpha_{k-1}$. Cut the word
$u_1\cdots u_{ak-1}$ at its descents into increasing runs. Each run has at most
$k$ vertices and there are at most $a$ runs, while there are $ak-1>k(a-1)$
vertices; so there are exactly $a$ runs, and only one place is missing from
the total capacity $ak$. The first run starts with colour $\alpha_1$, after
$\beta$, so it has at most $k-1$ vertices. Hence the first run is
$\alpha_1\cdots\alpha_{k-1}$ and every later run is
$\beta\alpha_1\cdots\alpha_{k-1}$; the colour word is
$A(\beta A)^{a-1}$ with $A=\alpha_1\cdots\alpha_{k-1}$. Now take the
permutation with $\alpha_1<\beta<\alpha_2<\dots<\alpha_{k-1}$. Each of the
$a-1$ factors $\alpha_{k-1}\beta\alpha_1$ contains the two descents
$\alpha_{k-1}\to\beta$ and $\beta\to\alpha_1$, so the path has at least
$2(a-1)>a-1$ descents and the direction is not good, a contradiction.
\end{proof}

\begin{lemma}\label{lem:rainbow}
For every vertex $v$, every proper $k$-colouring of $G-v$ and every $k$-cycle
$C$, the vertices of $C-v$ have distinct colours.
\end{lemma}

\begin{proof}
If $C$ avoids $v$, a good direction has one descent, so its $k$ colours are
distinct. If $C$ passes through $v$, the argument above shows that a good
direction needs $\alpha_1,\dots,\alpha_{k-1}$ distinct. If they were not
distinct, $C$ would be good for no permutation, contradicting
Corollary~\ref{cor:sat}.
\end{proof}


\paragraph{Successor digraphs.} Let $H$ be a graph with a proper
$k$-colouring $\varphi$ and let $\sigma$ be a cyclic permutation of the colours
(there are $(k-1)!$ of them). The digraph $D_\sigma(H,\varphi)$ has an arc
$s\to t$ whenever $st\in E(H)$ and $\varphi(t)=\sigma(\varphi(s))$. Every
directed cycle has length divisible by~$k$. If $R$ is a set of vertices closed
under out-arcs, replacing $\varphi$ by $\sigma\circ\varphi$ on $R$ gives
another proper colouring: an edge inside or outside $R$ stays proper, and an
edge $st$ with $s\in R$, $t\notin R$ that became monochromatic would satisfy
$\varphi(t)=\sigma(\varphi(s))$, an arc leaving~$R$.

\begin{lemma}\label{lem:edge}
Let $xy$ be an edge of $G$. Then at least $b$ $k$-cycles of $G$ contain $x$ and
avoid $y$.
\end{lemma}

\begin{proof}
Let $\varphi$ be a proper $k$-colouring of $G-xy$; then
$\varphi(x)=\varphi(y)$. Fix $\sigma$. If $y$ were not reachable from $x$ in
$D_\sigma(G-xy,\varphi)$, recolouring the set reachable from $x$ would give a
proper colouring of $G-xy$ in which $x$ and $y$ differ, hence a proper
$k$-colouring of $G$. So $y$ is reachable from $x$, and likewise $x$ from $y$.
A shortest closed directed walk through $x$ is a directed cycle through $x$,
and a cycle of $G$, since $D_\sigma$ has no loops and no directed $2$-cycles
($\sigma(c)\ne c$ and $\sigma^2(c)\ne c$ for $k\ge3$); by
Lemma~\ref{lem:short} it has length $k$, so its colours are distinct and it
avoids $y$, which has the colour of~$x$. The cyclic order of colours along the
cycle determines $\sigma$, and an undirected cycle arises from at most two
cyclic orders ($\sigma$ and $\sigma^{-1}$). Hence the $(k-1)!$ choices of
$\sigma$ give at least $b$ distinct $k$-cycles.
\end{proof}

\begin{proposition}\label{prop:degk}
If $G$ has a vertex of degree $k$, then $G\cong K_{k+1}$.
\end{proposition}

\begin{proof}
The cycles below are those constructed by Kim and
Picollelli~\cite[proof of Theorem~2.4 and Lemma~2.3]{kp}. Let $v$ have neighbours
$v_1,\dots,v_k$ and let $\varphi$ be a proper
$k$-colouring of $H=G-v$. All $k$ colours appear on $N(v)$ (otherwise $v$
could be coloured), so we may assume $\varphi(v_i)=i$.

\emph{Cycles through $v$.} For each $i$, colour $G_i=G-vv_i$ by $\varphi$
together with $\varphi(v)=i$. For each cyclic $\sigma$, the proof of
Lemma~\ref{lem:edge} gives a directed $k$-cycle through $v$ in
$D_\sigma(G_i,\cdot)$. Its vertices other than $v$ carry all colours except
$i$, so cycles obtained for different $i$ are different, and each $i$ gives at
least $b$ cycles. So at least $kb$ $k$-cycles pass through~$v$.

\emph{Cycles avoiding $v$.} Fix $\sigma$ and $i$. In $D_\sigma(G_i,\cdot)$
there is a directed path from $v_i$ to $v$. The vertex before $v$ has colour
$\sigma^{-1}(i)$, so it is $v_{\sigma^{-1}(i)}$, the only neighbour of $v$ of
that colour. Removing $v$ gives a directed path in $D_\sigma(H,\varphi)$ from
$v_i$ to $v_{\sigma^{-1}(i)}$. As $i$ varies these paths put
$v_1,\dots,v_k$ in one strong component of $D_\sigma(H,\varphi)$, which
therefore contains a directed cycle through each $v_i$. A directed cycle of
$D_\sigma$ has cyclic colour order $\sigma$, so cycles obtained from different
pairs $\{\sigma,\sigma^{-1}\}$ are different. Hence each pair
$\{\sigma,\sigma^{-1}\}$ yields a $k$-cycle avoiding~$v$, and at least $b$
$k$-cycles avoid~$v$.

\emph{Equality.} Since $N_k(G)=(k+1)b$, exactly $kb$ cycles pass through $v$
and exactly $b$ avoid it, so for each $\sigma$ the digraph $D_\sigma(H,\varphi)$
has exactly one undirected cycle among its directed cycles. That cycle passes
through all of $v_1,\dots,v_k$, has length $k$, and so equals $N(v)$ in the
colour order~$\sigma$; in particular $v_i$ is adjacent to
$v_{\sigma(i)}$. Every pair of distinct colours is consecutive in some
$\sigma$, so $N(v)$ is a clique and $\{v\}\cup N(v)$ spans $K_{k+1}$. By
criticality this subgraph is all of~$G$.
\end{proof}

\noindent Proposition~\ref{prop:degk} is used only for $k=4$ and $k=5$, and is
formalized only for these two values; for
$k\ge6$, Section~\ref{sec:six} gives more.

\begin{lemma}\label{lem:order}
Every vertex of $G$ lies on at least $b+1$ $k$-cycles.
\end{lemma}

\begin{proof}
Let $v$ be a vertex and $x$ a neighbour. Lemma~\ref{lem:edge} applied to the
edge $vx$ gives a $k$-cycle $C$ through $v$; let $y$ be a neighbour of $v$ on
$C$.
Lemma~\ref{lem:edge} applied to $vy$ gives $b$ further $k$-cycles through $v$
avoiding $y$, all different from~$C$.
\end{proof}

\section{The case $k=4$}\label{sec:four}

Let $k=4$, so $b=3$ and $B=15$. By Proposition~\ref{prop:degk} we may assume
that $\delta(G)\ge5$; then $n=|V(G)|\ge6$. By Lemmas~\ref{lem:short}
and~\ref{lem:order}, counting incidences between vertices and the $15$
four-cycles gives $4n\le60$, i.e.\ $n\le15$. For a pair $\{a,c\}$ of vertices let $t_{ac}$ be the
number of common neighbours. Every $4$-cycle is counted exactly twice by
choosing a pair of opposite vertices and two common neighbours, so
\begin{align*}
N_4(G)\ \ge\ \#\{4\text{-cycles}\}&=\tfrac12\sum_{\{a,c\}}\binom{t_{ac}}{2}
\ \ge\ \tfrac12\sum_{\{a,c\}}(t_{ac}-1)\\
&=\tfrac12\Bigl(\sum_w\binom{d(w)}{2}-\binom n2\Bigr)
\ \ge\ \tfrac12\Bigl(10n-\binom n2\Bigr)=\frac{n(21-n)}4 .
\end{align*}
On $6\le n\le15$ the right-hand side is at least $\frac{45}{2}$, so $G$ has at
least $23>15$ cycles of length $0$ modulo $4$, a contradiction. Hence
$\delta(G)=4$ and $G\cong K_5$.

\section{The case $k=5$}\label{sec:five}

In this section $k\ge5$, $G$ is $(k+1)$-critical with $N_k(G)=B$, and
$G\not\cong K_{k+1}$. We first prove three tools and then derive a
contradiction for $k=5$; the first tool is used again in
Section~\ref{sec:six}. We write
$n=|V(G)|$, $d_v$ for the degree of $v$, and $c_v$ for the number of
$k$-cycles through~$v$. By Lemma~\ref{lem:short} every cycle counted by
$N_k(G)$ is a $k$-cycle, so
\begin{equation}\label{eq:incidence}
\sum_{v} c_v = kB,
\end{equation}
which is $360$ for $k=5$. By
Proposition~\ref{prop:degk}, $\delta(G)\ge k+1$, hence $n\ge k+2$; by
Lemma~\ref{lem:order}, $c_v\ge b+1$ for every~$v$.

\subsection{Three tools}

\begin{lemma}[singleton neighbours]\label{lem:singleton}
Let $v$ be a vertex, $\varphi$ a proper $k$-colouring of $G-v$, and let $q$ be
the number of colours that occur exactly once on $N(v)$. Then
$q\ge 2k-d_v$ and $c_v\ge bq$. If $q\ge3$, the neighbours of $v$ with a
singleton colour, together with $v$, span a complete graph $K_{q+1}$.
\end{lemma}

\begin{proof}
All $k$ colours occur on $N(v)$, and a non-singleton colour occurs at least
twice, so $d_v\ge q+2(k-q)$. Let colour $i$ occur only at $x_i\in N(v)$.
Colouring $v$ with $i$ gives a proper $k$-colouring of $G-vx_i$. For each cyclic
permutation $\sigma$ of the colours, the proof of Lemma~\ref{lem:edge} gives a
directed $k$-cycle $C_i(\sigma)$ through $v$ in the successor digraph of this
colouring; it avoids $x_i$, and its vertices other than $v$ carry exactly the
colours different from~$i$. As in Lemma~\ref{lem:edge} each $i$ gives at least
$b$ distinct cycles, and cycles for different $i$ miss different colours, so
$c_v\ge bq$.

Let $j\ne i$ be another singleton colour. If $x_j\notin C_i(\sigma)$, then
colouring $v$ with $j$ and keeping $\varphi$ elsewhere is a proper colouring of
$G-x_j$, and $C_i(\sigma)$ contains $v$ and the vertex of $\varphi$-colour $j$
on it, which now have the same colour; this contradicts
Lemma~\ref{lem:rainbow}. So $C_i(\sigma)$ contains every other singleton
neighbour. If $q\ge3$, given singleton colours $j\ne l$ choose a third singleton
colour $i$ and $\sigma$ with $\sigma(j)=l$. The successor of $x_j$ on
$C_i(\sigma)$ has colour $l\ne i$, so it is not $v$; as $C_i(\sigma)-v$ is
rainbow and contains $x_l$, the successor is $x_l$. Hence $x_jx_l\in E(G)$.
\end{proof}

\begin{lemma}[outside neighbours of a $k$-cycle]\label{lem:attach}
Let $k=5$, let $C$ be a $5$-cycle and $y\notin V(C)$. Then $y$ has at most
$k-3=2$ neighbours on~$C$.
\end{lemma}

\begin{proof}
Since $n\ge k+2$ there is $w\notin V(C)\cup\{y\}$; let $\varphi$ be a proper
$k$-colouring of $G-w$. By Lemma~\ref{lem:rainbow}, $C$ carries all $k$ colours,
so exactly one vertex $s_0$ of $C$ has the colour of $y$; write
$C=s_0s_1\cdots s_{k-1}s_0$. Then $y\not\sim s_0$. If $y\sim s_j$ and
$y\sim s_{j+2}$ for some $1\le j\le k-3$, then
$y\,s_{j+2}s_{j+3}\cdots s_{k-1}s_0s_1\cdots s_j\,y$ is a $k$-cycle avoiding
$w$ that contains $y$ and $s_0$, contradicting Lemma~\ref{lem:rainbow}. So the
indices of the neighbours of $y$ form an independent set in the graph on
$\{1,\dots,k-1\}$ with edges $\{j,j+2\}$, which is the disjoint union of the
paths on the odd and on the even indices. For $k=5$ its independence number
is~$2$.
\end{proof}

\begin{lemma}[two double counts]\label{lem:double}
Let $k=5$. For every vertex $v$,
\[
b\,d_v\le (k-3)(B-c_v)\qquad\text{and}\qquad d_v(c_v-b)\ge 2c_v .
\]
\end{lemma}

\begin{proof}
For each $x\in N(v)$, Lemma~\ref{lem:edge} gives at least $b$ $k$-cycles
through $x$ avoiding $v$. There are $B-c_v$ $k$-cycles avoiding $v$, and by
Lemma~\ref{lem:attach} each contains at most $k-3$ neighbours of $v$. For the
second inequality, each $k$-cycle through $v$ contains at least two neighbours
of $v$, and each $x\in N(v)$ lies on at most $c_v-b$ of the cycles through $v$,
since at least $b$ of them avoid~$x$.
\end{proof}

\subsection{The case $k=5$}

Here $b=12$, $B=72$, $\sum_v c_v=360$, $\delta(G)\ge6$, $n\ge7$, and
Lemmas~\ref{lem:singleton} and~\ref{lem:double} read
\begin{equation}\label{eq:five-three}
c_v\ge12(10-d_v),\qquad c_v\le72-6d_v,\qquad
c_v\ge\Bigl\lceil\tfrac{12d_v}{d_v-2}\Bigr\rceil .
\end{equation}

\emph{Step 1: degrees $8$ and $9$ only, and $9\le n\le22$.} For $d_v=6$,
\eqref{eq:five-three} gives $c_v\ge48$ and $c_v\le36$; for $d_v=7$, $c_v\ge36$
and $c_v\le30$; for $d_v\ge10$, $c_v\le12$ and $c_v\ge13$. So
$d_v\in\{8,9\}$; if $d_v=8$ then $c_v=24$, and if $d_v=9$ then
$c_v\ge\lceil108/7\rceil=16$. Hence $16n\le360$, so $n\le22$, and $n\ge9$
since $\delta\ge8$.

\emph{Step 2: a budget for $5$-sets.} Let $S$ be the vertex set of a $5$-cycle,
$e$ the number of edges of $G[S]$, and $r$ the number of vertices outside $S$
with exactly two neighbours in $S$ (none has more, by Lemma~\ref{lem:attach}).
Then $\sum_{x\in S}d_x\le 2e+(n-5)+r$; using $c_x\ge12(10-d_x)$ on $S$ and
$c_x\ge13$ outside,
\begin{equation}\label{eq:five-budget}
\sum_x c_x\ \ge\ 595-24e+n-12r .
\end{equation}
In particular, if every outside vertex has at most one neighbour in $S$
then $\sum_x c_x\ge 12(50-n-15)+13(n-5)=355+n>360$, a contradiction.

\emph{Step 3: forbidden configurations on a $5$-cycle.} The common device is:
if $y\notin S$ has two neighbours $p,q\in S$, delete a vertex $w$ outside $S$
and outside the few further vertices involved (possible since $n\ge9$) and
colour $G-w$; $y$ has the colour of some $s\in S\setminus\{p,q\}$ ($S$ is
rainbow by Lemma~\ref{lem:rainbow}), and a path of length $3$ from $p$ to $q$
through $s$ inside $G[S]$ would close with $y$ into a $5$-cycle containing two
vertices of one colour. Using this,
\begin{itemize}
\item $G$ contains no $K_5$ and no $K_5-e$: in $K_5-e$ every such path exists,
so every outside vertex has at most one neighbour in $S$, contradicting Step 2.
\item No $5$-cycle has three chords ($e=8$). If the two non-edges are disjoint,
the same holds. If they share a vertex $z$ (so $S$ is a $K_4$ on $a,b,c,d$ plus
$z$ adjacent to $a,b$), every two-neighbour outside vertex has neighbourhood
$\{a,b\}$ or $\{c,d\}$ in $S$ and the colour of $z$; these vertices together
with $z$ are independent and have pairwise disjoint neighbourhoods outside $S$,
of size at least $4$, so $n\ge9+5r$ and $r\le2$; then \eqref{eq:five-budget}
with $e=8$ gives $\sum c_x\ge 403+9-24=388>360$.
\item $G$ has no gem (a path $abcd$ plus a vertex $z$ adjacent to all four).
Here the two-neighbour outside vertices split into two classes $R_a,R_d$, forced
to the colours of $a$ and $d$; the same disjointness gives
$|R_a|,|R_d|\le2$, so $r\le4$, and \eqref{eq:five-budget} with $e=7$ gives
$\sum c_x\ge 427+9-48=388>360$.
\end{itemize}
Hence every $5$-cycle has at most two chords, and two chords have no common
endpoint.

\emph{Step 4: a light edge on every $5$-cycle.} Every $5$-cycle $C$ has an edge
$uv$ such that $u,v$ have at most one common neighbour, and such a common
neighbour, if any, lies on $C$. If
$C$ has at most one chord and every edge of $C$ had a common neighbour off $C$,
these five vertices would be distinct (a vertex serving two edges would have
three neighbours on $C$, against Lemma~\ref{lem:attach}) and would form with $C$
a $10$-cycle, against Lemma~\ref{lem:short}. If $C=0\,1\,2\,3\,4$ has the two
chords $02$ and $13$, the edge $01$ works: its only common neighbour on $C$ is
$2$, and a common neighbour $w$ off $C$ would give a gem with centre $1$ on the
path $3\,2\,0\,w$.

\emph{Step 5: a Kempe-chain count.} Let $uv$ be such an edge and $\varphi$ a
proper $5$-colouring of $G-uv$, with $\varphi(u)=\varphi(v)=0$. For each other
colour $\beta$, $u$ and $v$ have $\beta$-neighbours and are joined by a shortest
$\{0,\beta\}$-path $P_\beta$ (otherwise a Kempe swap colours $G$). If $u,v$ have
no common $\beta$-neighbour then $|P_\beta|\ge6$ (length $4$ would give a
$5$-cycle through $uv$ that is not rainbow in $G-u$), so colour $\beta$ has at
least $a_\beta+b_\beta+1$ vertices, where $a_\beta,b_\beta$ count the
$\beta$-neighbours of $u,v$. For the colour $\beta_0$ missing on $C-u$, the path
$P_{\beta_0}$ together with $C-uv$ would form a $10$-cycle if it had length $6$,
so it has length at least $8$; hence colour $0$ has at least $5$ vertices and
$\beta_0$ at least $a_{\beta_0}+b_{\beta_0}+2$. If $u,v$ have a common
neighbour (necessarily on $C$, of a colour $\gamma\ne\beta_0$), colour $\gamma$
has at least $a_\gamma+b_\gamma-1$ vertices. Summing over the five colour
classes, with $\sum_{\beta\ne0}(a_\beta+b_\beta)=d_u+d_v-2$, gives
\[
n\ \ge\ d_u+d_v+6 .
\]

\emph{Step 6: contradiction.} By Step 1, $16\le d_u+d_v\le n-6\le16$. So
$n=22$ and $d_u=d_v=8$, whence $c_u=c_v=24$ while the other $20$ vertices have
$c_x\ge16$: $\sum_x c_x\ge24+24+20\cdot16=368>360$, contradicting
\eqref{eq:incidence}.


\section{Minimum degree when $k\ge6$}\label{sec:six}

Throughout this section $k\ge6$, $G$ is $(k+1)$-critical with $N_k(G)=B$, and
$G\not\cong K_{k+1}$. We prove:

\begin{proposition}\label{prop:mindeg}
Every vertex of $G$ has degree at least $2k-1$.
\end{proposition}

In this section we write $h=\frac{(k-1)!}{2}$ for the number called $b$ in
Section~\ref{sec:general}, as the letter $b$ will denote a vertex; so
$B=(k+1)h$. Fix a vertex $v$ and a proper $k$-colouring $\varphi$ of $G-v$. For a colour
$i$ let $F_i=\{x\in N(v):\varphi(x)=i\}$ and
$W_i=\{x\in V(G)-v:\varphi(x)=i\}$; note $F_i\subseteq W_i$. Every $F_i$ is
nonempty, since otherwise $v$ could be coloured~$i$. A colour $i$ is a
\emph{singleton colour} if $|F_i|=1$, and the vertex of $F_i$ is then its
\emph{centre}. Two facts are used throughout: by Lemma~\ref{lem:short} every
cycle counted by $N_k(G)$ has length exactly $k$, and by
Lemma~\ref{lem:rainbow} the vertices of $C-v$ have distinct colours for every
$k$-cycle~$C$ (and every choice of deleted vertex and colouring).

\subsection{Singleton colours}

\begin{lemma}\label{lem:retain}
Let $\alpha\ne\beta$ be singleton colours with centres $a$ and $b$. Then every
$k$-cycle through $v$ that avoids $a$ contains $b$ and contains no vertex
of~$W_\alpha$.
\end{lemma}

\begin{proof}
Colouring $v$ with $\alpha$ and every other vertex by $\varphi$ is a proper
colouring $\psi_a$ of $G-a$, since $a$ is the only neighbour of $v$ of colour
$\alpha$. Let $C$ be a $k$-cycle through $v$ avoiding $a$. Applying
Lemma~\ref{lem:rainbow} to the deleted vertex $a$ and the colouring $\psi_a$,
the $k$ vertices of $C$ have distinct $\psi_a$-colours. So $C$ contains no vertex
of $W_\alpha$, and it contains a vertex $z\ne v$ of colour~$\beta$. If $C$
avoided $b$, the same argument with the colouring $\psi_b$ of $G-b$ (colour $v$
with $\beta$) would show that $v$ is the only vertex of colour $\beta$ on $C$,
a contradiction.
\end{proof}

\begin{lemma}[donor paths]\label{lem:donor}
Let $\alpha\ne\beta$ be singleton colours with centres $a$ and $b$, and let
$\delta_1,\dots,\delta_s$ be distinct colours different from $\alpha$ and
$\beta$, where $s\ge0$ and $s+2\le k$. Then $G$ has a path
$v\,u_1\cdots u_s\,b$ avoiding $a$ with $\varphi(u_j)=\delta_j$ for all~$j$.
\end{lemma}

\begin{proof}
Colouring $v$ with $\alpha$ gives a proper colouring of $G-va$, in which $v$ and
$a$ have the same colour. Choose a cyclic permutation $\sigma$ of the colours
with $\sigma(\alpha)=\delta_1$, $\sigma(\delta_j)=\delta_{j+1}$ and
$\sigma(\delta_s)=\beta$ (with $\sigma(\alpha)=\beta$ if $s=0$). As in the
proof of Lemma~\ref{lem:edge}, the successor digraph of this colouring contains
a directed $k$-cycle through $v$; its colours are distinct, so it avoids $a$,
and by Lemma~\ref{lem:retain} it contains $b$, which is its only vertex of
colour~$\beta$. Following the cycle from $v$, the colours are
$\alpha,\delta_1,\dots,\delta_s,\beta$, so the first $s+2$ vertices form the
required path.
\end{proof}

\noindent For $s=1$ the lemma says that for every colour $\gamma\ne\alpha,\beta$
the centre $b$ has a neighbour in~$F_\gamma$.

\begin{lemma}[short spokes]\label{lem:spoke}
Let $\alpha\ne\beta$ be singleton colours with centres $a$ and $b$, and let $Q$
be a path from $a$ or from $b$ to $v$ of length $L$ with $2\le L\le k-1$. Then
the vertices of $Q-v$ have distinct colours.
\end{lemma}

\begin{proof}
By symmetry let $Q$ start at $b$. Suppose the claim fails, and let $T$ be the set of colours on $Q-v$; then $|T|\le L-1$ and
$\beta\in T$. Let $s=k-L-1\ge0$. Since $|T\cup\{\alpha\}|\le L$, there are at
least $k-L>s$ colours outside $T\cup\{\alpha\}$; choose $s$ of them. By
Lemma~\ref{lem:donor} there is a path $P$ from $v$ to $b$ of length $s+1$
avoiding $a$ whose inner vertices have these colours. The inner vertices of
$P$ and $Q$ have different colours, so $P$ followed by $Q$ is a cycle of length
$k$ through $v$, and it contains two vertices of $Q-v$ of the same colour,
contradicting Lemma~\ref{lem:rainbow}.
\end{proof}

\begin{lemma}\label{lem:twosingletons}
There are at most two singleton colours. Hence $d(v)\ge 2k-2$, and if
$d(v)=2k-2$ then there are exactly two singleton colours and $|F_i|=2$ for
every other colour~$i$.
\end{lemma}

\begin{proof}
The second statement follows from the first, since
$d(v)=\sum_i|F_i|\ge q+2(k-q)$ when $q$ colours are singleton colours.

Suppose that the set $J$ of centres has at least three elements. By
Lemma~\ref{lem:singleton}, $J\cup\{v\}$ is a clique. Call a colour that is not a
singleton colour \emph{plural}.

\emph{A spine.} Let $a\ne b$ be centres and $\gamma$ a plural colour. We find a
path from $a$ to $b$ of length $k-4$ avoiding $v$ and~$F_\gamma$. If
$|J|\ge k-3$, take a path through $k-3$ vertices of the clique $J$ from $a$ to
$b$. Otherwise let $s=k-3-|J|\ge1$; there are $k-|J|-1\ge s$ colours that are
neither the colours of $J$ nor $\gamma$, and we choose $s$ of them, say $D$.
Order $J$ as a path $a=w_0,w_1,\dots,w_{|J|-1}=b$, choose a third centre $h\ne
a,w_1$, and apply Lemma~\ref{lem:donor} to the centres $h$ (avoided) and $w_1$
(reached) with the colour sequence $\varphi(a),D$. This gives a path from $v$ to
$w_1$ whose first inner vertex is $a$, the only neighbour of $v$ of its colour,
and whose other inner vertices have the colours of $D$; removing $v$ and continuing along $w_1,\dots,w_{|J|-1}$ gives a
path of length $(s+1)+(|J|-2)=k-4$, since the colours of $D$ are not colours
of~$J$. No vertex of this path has colour~$\gamma$.

\emph{Hooks.} Let $\gamma$ be plural and let $p\ne p'$ in $F_\gamma$ be
adjacent to centres $a$ and $b$ respectively. If $a\ne b$, the cycle
$v\,p\,a\,(\text{spine})\,b\,p'\,v$ has length $k$ and contains two vertices of
colour~$\gamma$, contradicting Lemma~\ref{lem:rainbow}. By
Lemma~\ref{lem:donor} with $s=1$ every centre has a neighbour in $F_\gamma$. It
follows that there is a vertex $P_\gamma\in F_\gamma$ adjacent to every centre,
and no other vertex of $F_\gamma$ is adjacent to a centre.

\emph{The vertices $P_\gamma$ are pairwise adjacent.} Let $\gamma\ne\delta$ be
plural and let $a\ne b$ be centres with colours $\alpha,\beta$. By
Lemma~\ref{lem:donor} (avoiding $a$, reaching $b$, colours $\gamma,\delta$)
there is a path $v\,p\,z\,b$ with $\varphi(p)=\gamma$ and
$\varphi(z)=\delta$. If $p\ne P_\gamma$, the path $a\,P_\gamma\,b\,z\,p\,v$ of
length $5\le k-1$ contradicts Lemma~\ref{lem:spoke}. So $p=P_\gamma$. If $z\ne P_\delta$, the path
$b\,z\,P_\gamma\,a\,P_\delta\,v$ contradicts Lemma~\ref{lem:spoke} in the same
way. So $z=P_\delta$, and $P_\gamma P_\delta=pz$ is an edge.

Now the centres together with the vertices $P_\gamma$ for the plural colours
$\gamma$ are $k$ pairwise adjacent vertices, all adjacent to $v$; so $G$
contains $K_{k+1}$. If $G$ had a further vertex $w$, the graph $G-w$ would
contain $K_{k+1}$ and could not be $k$-coloured; so $G\cong K_{k+1}$, contrary
to our assumption.
\end{proof}

\subsection{Colour classes}

Let $c_x$ be the number of $k$-cycles through a vertex $x$. By
Lemma~\ref{lem:order}, $c_x\ge h+1$ for all~$x$.

\begin{lemma}\label{lem:classes}
$|W_i|\le k$ for every colour $i$. If $\alpha\ne\beta$ are singleton colours,
then $|W_\alpha|\le k-2$ and $|W_\beta|\le k-2$. In particular, if there are two
singleton colours, then $|V(G)|\le k^2-3$.
\end{lemma}

\begin{proof}
By Lemma~\ref{lem:rainbow} every $k$-cycle contains at most one vertex of
$W_i$, so $\sum_{x\in W_i}c_x\le B=(k+1)h$; if $|W_i|\ge k+1$ the left side
would be at least $(k+1)(h+1)$.

Let $a,b$ be the centres of $\alpha,\beta$. By Lemma~\ref{lem:edge} at least
$h$ $k$-cycles contain $v$ and avoid $b$, and by Lemma~\ref{lem:retain} they
all contain $a$; at least $h$ further $k$-cycles contain $a$ and avoid $v$. So
$c_a\ge2h$. At least $h$ $k$-cycles contain $v$ and avoid $a$, and by
Lemma~\ref{lem:retain} they contain no vertex of $W_\alpha$. Hence
$\sum_{x\in W_\alpha}c_x\le B-h$. If $|W_\alpha|\ge k-1$, the left side would
be at least $2h+(k-2)(h+1)=kh+k-2>kh=B-h$. Finally
$|V(G)|-1=\sum_i|W_i|\le 2(k-2)+(k-2)k=k^2-4$.
\end{proof}

\subsection{A vertex of degree $2k-2$}

We now assume $d(v)=2k-2$ and derive a contradiction, which proves
Proposition~\ref{prop:mindeg}. By Lemma~\ref{lem:twosingletons}, applied to
every vertex, every vertex has degree at least $2k-2$, there are exactly two singleton colours $\alpha,\beta$
with centres $a,b$, and every other colour $i$ has $|F_i|=2$. A vertex of
$N(v)-\{a,b\}$ adjacent to $a$ is an \emph{$a$-hook}, one adjacent to $b$ a
\emph{$b$-hook}. By Lemma~\ref{lem:donor} with $s=1$, every $F_i$
($i\ne\alpha,\beta$) contains an $a$-hook and a $b$-hook. By
Lemma~\ref{lem:classes}, $|V(G)|\le k^2-3$.

\emph{(i) $G$ contains no $K_k$.} Let $Q$ be a clique of size $k$. Since
$|V(G)|\ge d(v)+1=2k-1$, at least two vertices lie outside~$Q$. Suppose an
outside vertex $z$ has two neighbours $a_0,b_0\in Q$. Choose another outside
vertex $w$ and a proper $k$-colouring of $G-w$. The clique $Q$ uses all $k$
colours, so some $u\in Q$ has the colour of $z$, and $u\ne a_0,b_0$. Choose
$d_0\in Q-\{a_0,b_0,u\}$ and a path from $a_0$ to $b_0$ through all vertices of
$Q-d_0$; with $z$ it forms a $k$-cycle avoiding $w$ that contains $u$ and $z$,
contradicting Lemma~\ref{lem:rainbow}. So every outside vertex has at most one
neighbour in $Q$, and at most $|V(G)|-k\le k^2-k-3$ edges leave $Q$. But every
vertex of $Q$ has degree at least $2k-2$, hence at least $k-1$ neighbours
outside $Q$, so at least $k(k-1)$ edges leave~$Q$, a contradiction.

\emph{(ii) No vertex is both an $a$-hook and a $b$-hook.} Suppose $P\in F_\gamma$
is adjacent to $a$ and $b$, and let $\delta\ne\alpha,\beta,\gamma$. By
Lemma~\ref{lem:donor} (avoiding $a$, reaching $b$, colours $\gamma,\delta$)
there is a path $v\,p\,z\,b$ with $\varphi(p)=\gamma$, $\varphi(z)=\delta$. As
in the proof of Lemma~\ref{lem:twosingletons}, Lemma~\ref{lem:spoke} applied to
$a\,P\,b\,z\,p\,v$ gives $p=P$, and applied to $b\,z\,P\,a\,T\,v$, for any
$a$-hook $T\in F_\delta$, gives $z=T$. So the $a$-hook $T$ of colour $\delta$
is unique, is also a $b$-hook, and is adjacent to $P$. Applying this to each
such $T$ in place of $P$, the $a$-hooks of the $k-2$ colours other than
$\alpha,\beta$ are pairwise adjacent; together with $a$ and $v$ they form a
$K_k$, contradicting~(i).

\emph{(iii) No $a$-hook is adjacent to a $b$-hook.} Let $x$ be an $a$-hook and
$y$ a $b$-hook with $xy\in E(G)$; then $\varphi(x)\ne\varphi(y)$. Let $T$ be an
$a$-hook of colour $\varphi(y)$; by (ii), $T\ne y$. The path $b\,y\,x\,a\,T\,v$
has length $5\le k-1$ and contains $y$ and $T$, of the same colour,
contradicting Lemma~\ref{lem:spoke}.

Consequently, for every colour $i\ne\alpha,\beta$ we have $F_i=\{P_i,Q_i\}$,
where $P_i$ is adjacent to $a$ and not to $b$, $Q_i$ is adjacent to $b$ and not
to $a$, and no edge joins a vertex $P_i$ to a vertex~$Q_j$.

\emph{(iv) Locked colours.} For $i\ne\alpha,\beta$ let $K_i$ be the component
of the subgraph of $G-v$ induced by the colours $\alpha$ and $i$ that contains
$a$; it contains $P_i$. Call $i$ \emph{locked at $a$} if $Q_i\in K_i$. If $i$
is not locked at $a$, exchange the colours $\alpha$ and $i$ on $K_i$. This
gives a proper colouring $\varphi'$ of $G-v$ in which the colour classes on
$N(v)$ are $\{P_i\}$ for $\alpha$, $\{b\}$ for $\beta$, $\{a,Q_i\}$ for $i$,
and unchanged for the other colours. So $P_i$ and $b$ are the centres of the
two singleton colours of $\varphi'$, and Lemma~\ref{lem:donor} (with $s=1$)
for $\varphi'$ gives, for every $j\ne\alpha,\beta,i$, a neighbour of $P_i$ in
$F_j=\{P_j,Q_j\}$; it is not $Q_j$, so $P_iP_j\in E(G)$. Hence if
$P_iP_j\notin E(G)$, both $i$ and $j$ are locked at $a$. If all the $P_i$ were
pairwise adjacent, they would form with $a$ and $v$ a $K_k$; so by (i) at least
two colours are locked at $a$. In the same way (with $b$, $\beta$ and the
$Q_i$) at least two colours are \emph{locked at $b$}, that is, $P_j$ lies in
the component of the colours $\beta,j$ containing~$b$.

\emph{(v) Locked paths have length $k-1$.} Let $i$ be locked at $a$ and let
$L$ be a shortest path from $a$ to $Q_i$ in $K_i$, of length $x$. Its colours
alternate between $\alpha$ and $i$, so $x$ is odd, and $x\ge3$ since $a$ and
$Q_i$ are not adjacent. If $x\le k-2$, then $L$ followed by the edge $Q_iv$ is a
path from $a$ to $v$ of length at most $k-1$ whose at least four vertices other
than $v$ carry only two colours, contradicting Lemma~\ref{lem:spoke}. So
$x\ge k-1$. The path $L$ has $(x+1)/2$ vertices of colour $\alpha$, so
$(x+1)/2\le|W_\alpha|\le k-2$ by Lemma~\ref{lem:classes}, i.e.\ $x\le2k-5$.
Suppose $x\ge k$ and let $s=2k-x-3$, so that $2\le s\le k-3$. Choose $s$ colours
other than $\alpha,\beta,i$ and, by Lemma~\ref{lem:donor} with the roles of $a$
and $b$ exchanged, a path from $v$ to $a$ of length $s+1$ avoiding $b$ whose
inner vertices have these colours. Together with the path $L$, the edge $Q_ib$
and the edge $bv$ it forms a cycle of length $(s+1)+x+2=2k$, contradicting
Lemma~\ref{lem:short}. Hence $x=k-1$; since $x$ is odd, $k$ is even. The same
holds for a colour locked at~$b$.

\emph{(vi) A cycle of length $2k$.} Let $i$ be locked at $a$; as at least two
colours are locked at $b$, choose one, $j\ne i$. By (v) there are paths $L$ from
$a$ to $Q_i$ with colours $\alpha,i$ and $M$ from $b$ to $P_j$ with colours
$\beta,j$, both of length $k-1$. They are disjoint, as the four colours are
distinct, so $a\,L\,Q_i\,b\,M\,P_j\,a$ is a cycle of length $2k$, contradicting
Lemma~\ref{lem:short}. This proves Proposition~\ref{prop:mindeg}.

\section{A long cycle of length $0$ modulo $k$}\label{sec:long}

A \emph{rooted graph} $(H,x,y)$ is a graph $H$ with distinct vertices $x,y$; it
is \emph{$2$-connected} if $H+xy$ is $2$-connected. Paths $P_0,\dots,P_{r-1}$
are \emph{admissible} if $|P_0|\ge2$ and $|P_i|=|P_0|+id$ for all $i$, where
$d\in\{1,2\}$ ($|P|$ is the number of edges).

\begin{theorem}[{Gao, Huo, Liu and Ma~\cite[Theorem~3.1]{ghlm}}]\label{thm:ghlm}
Let $r\ge1$. If $(H,x,y)$ is a $2$-connected rooted graph in which every vertex
other than $x$ and $y$ has degree at least $r+1$, then $H$ has $r$ admissible
paths from $x$ to~$y$.
\end{theorem}

For $r=2$ this is \cite[Theorem~2.1]{ghlm}, which Gao, Huo, Liu and Ma derive
from a theorem of Fan~\cite[Theorem~2.5]{fan}. The formalization contains a
complete proof of Theorem~\ref{thm:ghlm}, including the case $r=2$; it
proceeds by induction on the number of vertices rather than on $|V|+|E|$ as
in~\cite{ghlm}.

\begin{lemma}\label{lem:progression}
Let $H$ be a connected graph of minimum degree at least $m\ge2$ and let
$2\le t\le m-1$. Then $H$ has cycles $C_0,\dots,C_{t-1}$ with $|C_i|=A+id$ for
some $A\ge3$ and $d\in\{1,2\}$.
\end{lemma}

\begin{proof}
This is the proof of~\cite[Theorem~1.3]{ghlm}. Choose a vertex $c$ and a
component $Z$ of $H-c$; among all sets $I\subseteq Z$ that are a component of
$H-c'$ for some vertex $c'$, choose one, $I$, of minimum size, with its vertex
$c'$. No edge leaves $I$ except to $c'$, and $H-I$ is connected (a path between
two vertices outside $I$ that enters $I$ must enter and leave through $c'$). If
a vertex $u\in I$ separated $H$, a component of $H-u$ not containing $c'$ would
miss $H-I$, hence be a proper subset of $I$, against minimality. It follows
that $H'=H[I\cup\{c'\}]$ is $2$-connected; $|I|\ge2$ because a single vertex
of $I$ would have degree at most~$1$. Every vertex of $I$ has all its neighbours in
$I\cup\{c'\}$, so its degree in $H'$ is at least $m\ge t+1$. Choose a neighbour
$y$ of $c'$ in $I$. Theorem~\ref{thm:ghlm} gives $t$ admissible paths from $c'$
to $y$ in $H'$; none of them uses the edge $c'y$, so adding this edge closes
each into a cycle, of length one more.
\end{proof}

\begin{proposition}\label{prop:long}
Let $k\ge4$ and let $G$ be a connected graph of minimum degree at least
$2k-1$. Then $G$ has a cycle whose length is a multiple of $k$ larger than~$k$.
\end{proposition}

\begin{proof}
\emph{$k$ odd.} Apply Lemma~\ref{lem:progression} to $G$ with $m=2k-1$ and
$t=2k-2$. We find an index $i<2k-2$ such that $k$ divides $A+id$ and
$A+id>k$. If $d=1$: when $A\le2k$ take $i=2k-A\le 2k-3$; otherwise take
$1\le i\le k$ with $k\mid A+i$. If $d=2$ and $A\le k$: take $i=(2k-A)/2$ if $A$
is even and $i=(3k-A)/2\le 2k-3$ if $A$ is odd. If $d=2$ and $A>k$: let
$1\le j\le k$ with $k\mid A+j$, and take $i=j/2$ if $j$ is even and $i=(j+k)/2$
if $j$ is odd.

\emph{$k$ even.} Among all partitions of $V(G)$ into two parts choose one with
the largest number of crossing edges, and let $H$ be the spanning subgraph of
the crossing edges. Moving one vertex to the other part shows
$d_H(u)\ge d_G(u)/2$, so $\delta(H)\ge k$. If $H$ were disconnected, moving one
of its components to the other side would keep all edges of $H$ crossing and
add an edge of $G$ joining that component to the rest, so $H$ is connected.
Apply Lemma~\ref{lem:progression} to $H$ with $m=k$ and $t=k-1\ge3$. As $H$ is
bipartite, all cycle lengths are even, so $A$ is even and $d=2$. Write $k=2q$
and $A=2c$ with $c\ge2$; among the $2q-1$ consecutive integers
$c,\dots,c+2q-2$ there is a multiple of $q$ larger than $q$ (take $2q$ if $c\le
2q$, and the next multiple of $q$ otherwise). Doubling gives a cycle of $H$,
hence of $G$, of the required length.
\end{proof}

\begin{remark}
For $2$-connected graphs with at least $2k+2$ vertices, Proposition~\ref{prop:long}
also follows from a theorem of Bai, Grzesik, Li and
Prorok~\cite[Theorem~1.6]{bglp}: for even $m\ge4$, every $2$-connected graph of
minimum degree at least $m-1$ with at least $m+2$ vertices has cycles of all
even lengths modulo $m$; with $m=2k$ this gives a cycle whose length is a
positive multiple of $2k$. We derive Proposition~\ref{prop:long} from
Theorem~\ref{thm:ghlm} because this is the route that is formalized.
\end{remark}

\section{Proof of Theorem~\ref{thm:main}}\label{sec:proof}

For $k=4$ and $k=5$ this is Sections~\ref{sec:four} and~\ref{sec:five}. Let
$k\ge6$, let $G$ be $(k+1)$-critical with $N_k(G)=B$, and suppose $G\not\cong
K_{k+1}$. The graph $G$ is connected: otherwise $k$-colourings of its
components, which are proper subgraphs, would combine into a $k$-colouring of
$G$. By Proposition~\ref{prop:mindeg}, $\delta(G)\ge2k-1$, so by
Proposition~\ref{prop:long} $G$ has a cycle whose length is a multiple of $k$
larger than $k$. This contradicts Lemma~\ref{lem:short}. \hfill$\square$

\medskip
Theorem~\ref{thm:main} is the uniqueness statement asked for in~\cite{cs}; that
$K_{k+1}$ is $(k+1)$-critical and has exactly $B$ cycles of length divisible by
$k$ is the observation recalled in the introduction.

\section{Formal verification}\label{sec:lean}

Theorem~\ref{thm:main} is formally verified in Lean~4 (version 4.34.1) with
Mathlib~\cite{mathlib} (version v4.34.1, commit \texttt{d13f23b7}). The $716$
files are in version~\repoversion{} of the repository~\cite{repo}, directory
\texttt{lean/\allowbreak Research/\allowbreak TaskH/}.

\paragraph{The statements.}
The statement file \texttt{Challenge.lean} contains only definitions and
statements, no proofs. In ASCII notation (implicit arguments omitted):
{\footnotesize
\begin{verbatim}
def IsCycleEdgeSet (G : SimpleGraph V) (s : Finset (Sym2 V)) : Prop :=
  exists (v : V) (p : G.Walk v v), p.IsCycle /\ p.edges.toFinset = s
def ModCycleEdgeSets (k : Nat) (G : SimpleGraph V) : Set (Finset (Sym2 V)) :=
  {s | IsCycleEdgeSet G s /\ k | s.card}
def modCycleCount (k : Nat) (G : SimpleGraph V) : Nat := (ModCycleEdgeSets k G).ncard
def IsCriticalPlus (k : Nat) (G : SimpleGraph V) : Prop :=
  G.chromaticNumber = (k + 1 : Nat) /\ forall S : G.Subgraph, S != top -> S.coe.Colorable k
def bound (k : Nat) : Nat := (k + 1).factorial / (2 * k)
def EqualityCase (k : Nat) : Prop :=
  forall (V : Type u) [Fintype V] (G : SimpleGraph V),
    IsCriticalPlus k G -> modCycleCount k G = bound k ->
      Nonempty (G =g (top : SimpleGraph (Fin (k + 1))))
def Question2 : Prop := forall k : Nat, 4 <= k -> EqualityCase k
\end{verbatim}}
\noindent (In the file the conclusion is written with an auxiliary definition
\texttt{IsCompleteOfOrder}, \texttt{G =g H} is a graph isomorphism, and
\texttt{top} is the complete graph.) A cycle is the edge set of a closed walk
that is a cycle in Mathlib's sense, so neither its starting vertex nor its
direction is counted. Criticality quantifies over all proper subgraphs,
including those obtained by deleting a single edge. The division in
\texttt{bound} is exact: the files prove $2k\cdot\texttt{bound}\,k=k!\,(k+1)$
for $k\ge3$. \texttt{Question2}, the files' name for Question~8
of~\cite{cs}, is Theorem~\ref{thm:main}. The files
\texttt{CaseFive.lean}, \texttt{CaseSix.lean} and \texttt{CaseSeven.lean}
define \texttt{EqualityCase 5}, \texttt{6} and \texttt{7}, which the first
version proved separately.

\paragraph{Review and freezing.}
\texttt{Challenge.lean}, including \texttt{Question2}, was reviewed against the
question by a Codex agent that had not written it and frozen, by recording its
SHA-256 hash, on 13 November 2026 at 18:42~UTC, before any proof file was
written; it was not changed afterwards. The files for $k=5,6,7$ were frozen in
the same way on 13 and 14 November.

\paragraph{The proofs.}
The theorem
\texttt{TaskH.\allowbreak original\_\allowbreak question\_\allowbreak two\_\allowbreak all\_\allowbreak k} in
\texttt{Original\allowbreak QuestionTwo\allowbreak Assembly.lean} proves \texttt{Question2}: it uses
\texttt{case\_\allowbreak four} and \texttt{case\_\allowbreak five} for $k=4,5$, and for $k\ge6$
combines \texttt{equality\_\allowbreak shortCycles\_\allowbreak from\_\allowbreak critical\_\allowbreak count}
(Lemma~\ref{lem:short}),
\texttt{critical\_\allowbreak original\_\allowbreak connected},
\texttt{equality\_\allowbreak original\_\allowbreak minimum\_\allowbreak degree\_\allowbreak two\_\allowbreak mul\_\allowbreak sub\_\allowbreak one}
(Proposition~\ref{prop:mindeg}) and
\texttt{connected\_\allowbreak high\_\allowbreak degree\_\allowbreak original\_\allowbreak long\_\allowbreak zero\_\allowbreak cycle}
(Proposition~\ref{prop:long}). The last uses
\texttt{gao\_\allowbreak original\_\allowbreak rooted\_\allowbreak path\_\allowbreak claim}, a proof of
Theorem~\ref{thm:ghlm}, and \texttt{fan\_\allowbreak rooted\_\allowbreak two\_\allowbreak paths}, a proof of its
case $r=2$ \cite[Theorem~2.1]{ghlm}; neither is assumed. The good-cycle lemmas of Section~\ref{sec:good} are
reproved, not assumed. The test file \texttt{OriginalQuestionTwoTests.lean}
restates the theorem with the frozen type and proves that the frozen
counterexample statement is false, and the files
\texttt{ExactTests.lean}, \ldots, \texttt{ExactTestsSeven.lean} check
$K_{k+1}$ for $k=4,\dots,7$: it is $(k+1)$-critical with exactly
\texttt{bound k} counted cycles ($15$, $72$, $420$ and $2880$). That $K_{k+1}$ is
$(k+1)$-critical is proved for every $k$ (\texttt{complete\_\allowbreak isCriticalPlus}),
but its cycle count is formalized only for $k\le7$. The axioms of
\texttt{original\_\allowbreak question\_\allowbreak two\_\allowbreak all\_\allowbreak k}, printed by the files, are exactly
\texttt{propext}, \texttt{Classical.choice} and \texttt{Quot.sound}. No file
uses \texttt{sorry}, \texttt{native\_\allowbreak decide}, \texttt{bv\_\allowbreak decide} or an
\texttt{axiom} declaration; \texttt{decide +kernel} is used, which is checked
by the kernel and adds no axiom. Comments in some files (for example
\texttt{Original\allowbreak QuestionTwo\allowbreak Assembly.lean},
\texttt{General/\allowbreak OriginalStrongerMinimumDegree.lean} and
\texttt{CaseSix.lean}) were written before the reviews and checks below and
call the files candidates, unchecked, or awaiting review; they were not
changed, so as not to change the checked files.

\paragraph{Checks.}
The $717$ files of Codex's accepted package (the $716$ files of the repository
and a statement file for $k=8$ that nothing imports and that is not included)
were compiled module by module in two new empty output directories on the
author's Mac, with the pinned toolchain and Mathlib, on 17 November 2026
(04:08--07:14~UTC); every module was replayed with \texttt{leanchecker}. All
$2868$ commands ended with exit code~$0$ and printed nothing except the
requested axiom lists, and the \texttt{.olean} files of the two directories are
byte for byte identical. In the whole-project rebuild of version~1.22.0 of the repository (17--18 November 2026, on the author's Windows computer), the $716$ files were compiled and replayed with \texttt{leanchecker} again, together with the rest of the project; runs that failed for lack of memory or file access, without a Lean verdict, were run again and passed (details in the repository). The formalization has not been checked
with comparator or with an independent kernel.

\section*{How this paper was produced}
This work was carried out from 13 to 18 November 2026 (UTC) in a
workflow directed by the author, using OpenAI Codex with the model GPT-6.1 Sol
and Anthropic Claude Code with the model Claude Opus~5.5. Claude Code's
scouting agents proposed the question; the author chose it, and Claude Code
assigned it to Codex. Codex agents found the proofs, wrote the statement files,
the mathematical manuscripts and the Lean files, and ran the checks. Each
statement file was reviewed before it was frozen by a Codex agent (same model)
that had not written it. The cases $k=4,\dots,7$ were reviewed after their
checks by Codex agents that were not among the recorded authors of those proofs
(details in version~1.21.0 of the repository, including that for $k=7$ the
records do not identify the authors of $40$ files). The proof for all $k\ge4$
was reviewed after its checks by an independent Claude Code agent (model Claude
Opus~5.5) that had written none of the files. It reread the raw logs of the
checks, read the statement file and the main proof files, read Codex's
written proof of Proposition~\ref{prop:mindeg} and the parts of its written
proof of Proposition~\ref{prop:long} that derive it from
Theorem~\ref{thm:ghlm}, and approved; it did not read the Lean proofs of
Theorem~\ref{thm:ghlm} and of its case $r=2$ line by line, nor the
deeper counting lemmas, whose correctness rests on the Lean check. Claude Code
searched the literature, wrote this text from Codex's manuscripts and the Lean
files, and rebuilt the whole Lean project of the repository; an independent
Claude Code agent (model Claude Opus~5.5) that had written none of the files reviewed this paper, and another reviewed the release; their corrections were applied. Throughout, ``we'' refers to this workflow;
the reading and searching described in the paper were done by the agents. The
author is not a professional mathematician. The author has read this paper and, using a side-by-side table, compared its Lean statements with Theorem~\ref{thm:main}. No human
mathematician has reviewed this paper. The correctness of the result rests on
the Lean proofs and on the agreement of the Lean statements with the question
of~\cite{cs}; the written proofs are given for the reader. AI systems are not
authors of this paper; the author takes full responsibility for its content.

\medskip
\noindent Thanks are due to Ilkyoo Choi and Kamil Sieniawski for the question.

\begin{thebibliography}{9}
\bibitem{bglp} Y.~Bai, A.~Grzesik, B.~Li and M.~Prorok, Cycle lengths in graphs
of given minimum degree, \emph{J. Combin. Theory Ser. B} 180 (2026) 111--150;
arXiv:2511.03085.
\bibitem{bmmn} R.~C. Brewster, S.~McGuinness, B.~Moore and J.~A. Noel,
A dichotomy theorem for circular colouring reconfiguration,
\emph{Theoret. Comput. Sci.} 639 (2016) 1--13.
\bibitem{cs} I.~Choi and K.~Sieniawski, The minimum number of cycles of length
$0$ modulo $k$ in graphs with no proper $k$-coloring, arXiv:2611.13580v1, 2026.
\bibitem{repo} H.~Dong, \emph{ai4math-results}, repository,
\url{https://github.com/Horace-Maxwell/ai4math-results}, version~\repoversion{}
(2026); archived at Zenodo under the concept DOI
\url{https://doi.org/10.5281/zenodo.22970254}.
\bibitem{fan} G.~Fan, Distribution of cycle lengths in graphs, \emph{J. Combin.
Theory Ser. B} 84 (2002) 187--202.
\bibitem{ghlm} J.~Gao, Q.~Huo, C.-H. Liu and J.~Ma, A unified proof of
conjectures on cycle lengths in graphs, \emph{Int. Math. Res. Not. IMRN}
2022 (2022), no.~10, 7615--7653; arXiv:1904.08126.
\bibitem{kp} S.~Kim and M.~E. Picollelli, $4$-chromatic graphs have at least
four cycles of length $0 \bmod 3$, \emph{Electron. J. Combin.} 31(2) (2024)
Paper No.~2.58; arXiv:2312.05945.
\bibitem{mathlib}
The mathlib Community, The Lean mathematical library, in \emph{Proc.\ 9th ACM
SIGPLAN Int.\ Conf.\ Certified Programs and Proofs} (CPP '20), ACM, 2020,
\url{https://doi.org/10.1145/3372885.3373824}.
\end{thebibliography}
\end{document}
